% Abstract (owner's text, verbatim). Prose, unlike the other section files.
\begin{abstract}
  A model can improve itself by splitting hard problems into easier ones it can
  already solve, combining their answers, and then training on the resulting
  data. \citet{shin2026compositional} showed that such training rounds can help,
  do nothing, or harm the model, and diagnosed the failures after training, using
  correct answers. We ask the complementary question: can the outcome of a
  training round be predicted before training, without access to correct
  answers, and at a lower cost than running the round itself? We view round $i$
  of the self-improvement loop as a teacher (the model $M_{i-1}$ applying the
  split-and-combine procedure, which produces the dataset $D_i$) and a student
  ($M_{i-1}$ answering directly). In this view, training on $D_i$ can help only
  if the teacher's answers differ from the student's, are correct, and include
  information the student lacks. The last two conditions restate failure modes
  identified by Shin et al.; the first is new. We propose a check for each,
  performed before training: comparing the teacher's answers with the student's,
  comparing different ways of splitting the same problem, and examining what the
  combining step is able to produce. None uses correct answers or requires
  knowing which problems the procedure gets wrong. Reproducing Shin et al.'s
  setup with a 14M-parameter Transformer on two tasks, and fixing all thresholds
  before the final evaluation, we find that our checks predict a round's outcome
  in 96.5\% of 258 cases and flag every harmful round, at about a tenth of the
  cost of the round itself. Starting the round and peeking at early progress is
  no more accurate, costs more than half a round, and requires correct answers.
  These results suggest that a recursively self-improving model need not run a
  round to learn its outcome: from its own answers alone, it could tell when
  another round will help, when it will do harm, and when no further round can
  help.



\end{abstract}
